For an integer $n\ge 4$, ask whether there exists a set of points
\[S=\{p_1,\dots,p_n\}\subset\mathbb{R}^2\]
such that:
\begin{enumerate}
  \item \textbf{Integral distances:} for all $1\le i<j\le n$, $d(p_i,p_j)\in\mathbb{Z}_{>0}$.
  \item \textbf{No three collinear:} for all distinct $i,j,k$, the points $p_i,p_j,p_k$ do not lie on a single line.
  \item \textbf{No four concyclic:} for all distinct $i,j,k,\ell$, the points $p_i,p_j,p_k,p_\ell$ do not lie on a single circle.
\end{enumerate}
The natural strengthened interpretation (and the one that makes this a genuine ``Erd\H{o}s problem'') is:
\begin{quote}
\emph{Do such configurations exist for arbitrarily large $n$ (equivalently, for every $n$ beyond some point)?}
\end{quote}
Indeed, if a configuration exists for some $N$, then any subset of size $n\le N$ is again in general position and still has all distances integral.
